For a $k$-critical graph $G$ on $n$ vertices, let the edge-bipartization cost be the fewest edges whose removal makes $G$ bipartite. Must its minimum $f_k(n)$ tend to infinity, and in particular must $f_4(n)\gg\log n$?
The answer is no. For every fixed $k\ge4$ and every $r\ge1$, the construction below is $k$-critical, has $r(k-1)+1$ vertices, and has bipartization cost at most $\binom{k-1}{2}$. Thus even an unbounded sequence of orders disproves both proposed growth conclusions.
The generalized Mycielski construction is classical. R\"odl and Tuza proved the sharper eventual identity $f_k(n)=\binom{k-1}{2}$; the proof here establishes the counterexamples directly and does not claim their sharp lower bound or their result at every sufficiently large order.
